% А1.1 — Формулы сокращённого умножения; вынесение общего множителя и группировка
% Старая методичка: стр. 49 (старое задание 6)
% Источники: шпаргалка «Алгебра», методичка «Деление уголком» (§ 2, подсказки).
% Стиль — стиль-изложения.md

\begin{tema}{А1.1}
\begin{pfTablica}[bez shapki]{rlrl}
\pfPodshapka{Квадраты и кубы суммы и разности}\\
\nomer1 & $(a+b)^2=a^2+2ab+b^2$ & \nomer3 & $(a+b)^3=a^3+3a^2b+3ab^2+b^3$\\
\nomer2 & $(a-b)^2=a^2-2ab+b^2$ & \nomer4 & $(a-b)^3=a^3-3a^2b+3ab^2-b^3$\\
\pfPodshapka{Разность квадратов, сумма и разность кубов}\\
\nomer5 & $a^2-b^2=(a-b)(a+b)$ & \nomer6 & $a^3-b^3=(a-b)(a^2+ab+b^2)$\\
 & & \nomer7 & $a^3+b^3=(a+b)(a^2-ab+b^2)$\\
\pfPodshapka{Квадрат трёхчлена}\\
\nomer8 & \Block[l]{1-3}{$(a+b+c)^2=a^2+b^2+c^2+2ab+2bc+2ac$}\\
\end{pfTablica}

Формулы работают в обе стороны: слева направо — чтобы раскрыть скобки,
справа налево — чтобы разложить выражение на множители.

\opr{Вынесение общего множителя:} $ab+ac=a(b+c)$.

\opr{Группировка.} Слагаемые объединяют в группы так, чтобы в каждой появился
одинаковый множитель, и выносят его: $ax+ay+bx+by=a(x+y)+b(x+y)=(a+b)(x+y)$.

\begin{primer}
Решите уравнение $x^3-4x^2-x+4=0$.
Решение: сгруппируем слагаемые:
$x^2(x-4)-(x-4)=(x-4)(x^2-1)=(x-4)(x-1)(x+1)=0$ — группировка, затем формула (5).
Ответ: $-1;\ 1;\ 4$.
\end{primer}
\end{tema}
